\chapter{轨迹规划}

轨迹规划的任务是生成一个随时间变化的位姿函数 $\theta(t)$，使得机器人能够平滑地从起点移动到终点。轨迹可以看作是“路径”与“时间规律”的结合。

\section{ 路径与轨迹的关系}
路径（Path）是构型空间中曲线的几何表示，不含时间信息。轨迹（Trajectory）则是参数化的路径：
\begin{equation}
    \theta(t) = s(t) \cdot \text{Path}
\end{equation}
其中 $s(t) \in [0, 1]$ 是时间标度函数（Time Scaling），描述了机器人在路径上的进度。

\section{点到点轨迹 (Point-to-Point Trajectories)}
对于给定的起点 $\theta_{start}$ 和终点 $\theta_{end}$，通常使用多项式插值来保证运动的平滑性。



\begin{mybox}{三次多项式轨迹 (Cubic Polynomials)}
若给定起止位置和速度（通常速度为 0），可使用 $s(t) = a_0 + a_1 t + a_2 t^2 + a_3 t^3$：
\begin{itemize}
    \item 满足 4 个约束：$s(0), s(T), \dot{s}(0), \dot{s}(T)$。
    \item 优点：计算简单，位置和速度连续。
    \item 缺点：加速度在起止时刻可能发生跳变。
\end{itemize}
\end{mybox}

\begin{mybox}{五次多项式轨迹 (Quintic Polynomials)}
若需保证加速度连续（即零加加速度），需使用 $s(t) = a_0 + a_1 t + \dots + a_5 t^5$：
\begin{itemize}
    \item 满足 6 个约束：包括起止时刻的位置、速度和加速度。
    \item 在高性能机器人控制中，为了减小机械振动，通常首选五次多项式。
\end{itemize}
\end{mybox}

\section{梯形速度曲线 (Trapezoidal Motion Profiles)}
在工程中，为了充分利用电机的额定转速和加速度，常采用梯形速度曲线（也称为恒定加速度轨迹）。



\begin{defbox}{梯形速度特征}
轨迹被分为三个阶段：
\begin{enumerate}
    \item \textbf{匀加速段}：速度线性上升，加速度为常数。
    \item \textbf{匀速段}：保持最大允许速度运行。
    \item \textbf{匀减速段}：速度线性下降至零。
\end{enumerate}
若最大速度不足以达到，则曲线退化为三角形速度曲线。
\end{defbox}

\section{多点路径平滑 (Multi-point Trajectories)}
当机器人需要经过一系列中间点（Via-points）时，直接使用直线段连接会导致在中间点处速度不连续。

\begin{thoughtbox}{关于“混合”段的思考}
为了在中间点处平滑过渡，通常引入“混合区域”（Blending regions）。在接近中间点时，提前开始向下一个目标点加速。这种方法虽然不会精确经过中间点，但能保证路径的连续性和执行效率。
\end{thoughtbox}



\section{任务空间轨迹规划}
在任务空间（笛卡尔空间）规划时，不仅要插值位置 $p(t)$，还要插值旋转矩阵 $R(t)$。

\begin{formulabox}{旋转插值 (SLERP)}
简单的欧拉角插值会导致奇异性，推荐使用轴角法或球面线性插值（SLERP）：
\begin{equation}
    R(t) = R_{start} \exp([\hat{\omega}]\theta t)
\end{equation}
这保证了末端执行器绕固定轴以恒定角速度旋转。
\end{formulabox}

\clearpage